\documentclass{oxmathproblems}

\printanswers

\course{Understanding Analysis, 2nd Edition}
\sheettitle{Chapter 6 - Section 6.4 - Series of Functions} 

\begin{document}
\begin{questions}

%------------------------- Problem 1 -------------------------
\miquestion Supply the details for the proof of the Weierstrass M-Test (Corollary 6.4.5).
\begin{solution}
  We have for every fixed $x$
  \[0 \leq |f_{n}(x)| \leq M_{n}\]
  and $\sum_{n=1}^{\infty}M_{n}$ converges. This satisfies the conditions for the Comparison Test, so $\sum_{n=1}^{\infty}|f_{n}(x)|$ converges. 

  Then by the Absolute Convergence Test, $\sum_{n=1}^{\infty}|f_{n}(x)|$ converges implies $\sum_{n=1}^{\infty}f_{n}(x)$ converges.

  It remains to prove that the convergence is uniform. $\sum_{n=1}^{\infty}M_{n}$ converges so it satisfies the Cauchy Criterion for Series (Theorem 2.7.2):
  for any $\epsilon > 0$, there exists $N$ such that for all $n > m \geq N$:
  \[M_{m+1} + M_{m+2} + ... + M_{n} < \epsilon\]
  Reusing the same $N$, we can prove that $\sum_{n=1}^{\infty}f_{n}(x)$ satisfies the Cauchy Criterion for Uniform Convergence of Series (Theorem 6.4.4). Because
  $|f_{n}(x)| \leq M_{n}$, then
  \begin{align*}
    |f_{m+1}(x) + f_{m+2}(x) + ... + f_{n}(x)| &\leq |f_{m+1}(x)| + |f_{m+2}(x)| + ... + |f_{n}(x)|\\
    &\leq M_{m+1} + M_{m+2} + ... + M_{n}\\
    &< \epsilon 
  \end{align*}
  Since the same $N$ works for any $x$ (because the bound $M_{n}$ is independent of $x$), the convergence is uniform.
\end{solution}

%------------------------- Problem 2 -------------------------
\miquestion Short proof or counterexample.
\begin{enumerate}[label=(\alph*)]
  \item If $\sum_{n=1}^{\infty}g_{n}$ converges uniformly, then $(g_{n})$ converges uniformly to zero.
  \item If $0 \leq f_{n}(x) \leq g_{n}(x)$ and $\sum_{n=1}^{\infty}g_{n}$ converges uniformly, then $\sum_{n=1}^{\infty}f_{n}$ converges uniformly.
  \item If $\sum_{n=1}^{\infty}f_{n}$ converges uniformly on $A$, then there exist constants $M_{n}$ such that $|f_{n}(x)| \leq M_{n}$ for all $x \in A$
  and $\sum_{n=1}^{\infty}M_{n}$ converges.
\end{enumerate}

\begin{solution}
  \begin{enumerate}[label=(\alph*)]
    \item True. This is analogous with the $n$th-term test for series of numbers (Theorem 2.7.3).
    
    $\sum_{n=1}^{\infty}g_{n}$ converges uniformly so there exists $N$ such that for all $n > m > N$ and for all $x$
    \[|g_{m+1}(x) + ... + g_{n}(x)| < \epsilon\]
    The same $N$ works to prove $(g_{n})$ converges uniformly. Set $n=m+1$.

    Alternatively, without using the Cauchy Criterion, let $s_{k} = \sum_{n=1}^{k}g_{n}(x)$. Then
    \[\lim_{k \to \infty}s_{k}-s_{k-1}=\lim_{k \to \infty}s_{k} - \lim_{k \to \infty}s_{k-1}=0\]
    And note that $s_{k}-s_{k-1}=g_{k}$, so $\lim_{k \to \infty}s_{k}-s_{k-1} = \lim_{k \to \infty}g_{k}(x)=0$.

    The important intuition here is that uniform convergence allows us to reuse many results from series of numbers,
    because any bound works uniformly so we may treat $g_{n}(x)$ for all $x$ like a single entity.
  
    \item True. This is analogous to the Comparison Test for series of numbers (Theorem 2.7.4).
    
    Like part (a), but we can get rid of the absolute values because it's all non-negative
    \[g_{m+1}(x) + ... + g_{n}(x) < \epsilon\]
    then reusing the same $N$ and the problem's assumption that $0 \leq f_{n}(x) \leq g_{n}(x)$
    \[f_{m+1}(x) + ... + f_{n}(x) \leq g_{m+1}(x) + ... + g_{n}(x) < \epsilon\]

    \item False. This is analogous to conditional convergence for series of numbers (Theorem 2.7.8).
    
    Take $f_{n}(x)=\frac{(-1)^{n+1}}{n}$ for all $x$.
    Then $\sum_{n=1}^{\infty}f_{n}$ is the alternating harmonic series for all $x$ so it converges uniformly.

    But $M_{n} \geq \frac{1}{n}$ so $\sum_{n=1}^{\infty}M_{n}$ diverges by comparison test with $\sum_{n=1}^{\infty}\frac{1}{n}$.
  \end{enumerate}
\end{solution}

%------------------------- Problem 3 -------------------------
\miquestion
\begin{enumerate}[label=(\alph*)]
  \item Show that
  \[g(x) = \sum_{n=0}^{\infty}\frac{\cos(2^{n}x)}{2^n}\]
  is continuous on all of \textbf{R}.
  \item The function $g$ was cited in Section 5.4 as an example of a continuous nowhere differentiable function. What happens if we try to use Theorem 6.4.3 to
  explore whether $g$ is differentiable?
\end{enumerate}

\begin{solution}
  \begin{enumerate}[label=(\alph*)]
    \item We have $|\cos(2^{n}x)| \leq 1$ for any $n$ and $x$. Then
    \[\left|\frac{\cos(2^{n}x)}{2^n}\right| \leq \frac{1}{2^n}\]
    for all $x$ and $\sum_{n=0}^{\infty}\frac{1}{2^n} = 2$. This satisfies the Weierstrass M-test, so $\sum_{n=0}^{\infty}\frac{\cos(2^{n}x)}{2^n}$ converges uniformly.

    By Theorem 6.4.2 (Term-by-term Continuity Theorem), since $\frac{\cos(2^{n}x)}{2^n}$ is continuous for any $n$ and $\sum_{n=0}^{\infty}\frac{\cos(2^{n}x)}{2^n}$
    converges uniformly, then $g$ is continuous.

    \item The series of derivative is
    \[\sum_{n=0}^{\infty}-\sin(2^{n}x)\]
    This fails to converge uniformly so we can't apply Theorem 6.4.3. For example, take $x=\pi/3$, then the terms are
    \[
      -\sin\left(\frac{2^{n}\pi}{3}\right)=\begin{cases}
        -\frac{\sqrt{3}}{2}, & n=0\\
        (-1)^{n}\frac{\sqrt{3}}{2}, & n \geq 1\\
      \end{cases}
    \]
    Since the terms fail to approach 0, the series fails to converge.

    This does not prove that $g$ is not differentiable; it only means we can't use Theorem 6.4.3 to conclude differentiability.
  \end{enumerate}
\end{solution}

%------------------------- Problem 4 -------------------------
\miquestion Define
\[g(x)=\sum_{n=0}^{\infty}\frac{x^{2n}}{(1+x^{2n})}\]
Find the values of $x$ where the series converge and show that we get a continuous function on this set.

\begin{solution}
  We have
  \[\frac{x^{2n}}{(1+x^{2n})} = \frac{x^{2n}+1-1}{(1+x^{2n})}=1-\frac{1}{1+x^{2n}}\]

  For $x \geq 1$ or $x \leq -1$, the terms converge to 1, so the series diverges.

  For $x \in (-1, 1)$ the series converges but not uniformly so we can't directly apply Theorem 6.4.2. But we can use the same trick as problem 5(b).
  Choose $r$ such that $|x| < r < 1$, then on $[-r, r]$
  \[0 \leq \frac{x^{2n}}{1+x^{2n}} \leq x^{2n} \leq r^{2n}\]
  Since $\sum_{n=0}^{\infty}r^{2n}$ is a convergent geometric series, we may apply the Weierstrass M-Test to conclude that $\sum_{n=0}^{\infty}\frac{x^{2n}}{(1+x^{2n})}$
  converges uniformly on $[-r, r]$.

  Since each term of the series is continuous, we may apply the Term-by-term Continuity Theorem to conclude that $g$ is continuous on $[-r, r]$ hence continuous on $x$.
  Since this argument works for any $x \in (-1, 1)$, $g$ is continuous on $(-1, 1)$.

\end{solution}

%------------------------- Problem 5 -------------------------
\miquestion 
\begin{enumerate}[label=(\alph*)]
  \item Prove that
  \[h(x) = \sum_{n=1}^{\infty}\frac{x^n}{n^2} = x + \frac{x^2}{4} + \frac{x^3}{9} + ...\]
  is continuous on $[-1, 1]$.

  \item The series
  \[f(x) = \sum_{n=1}^{\infty}\frac{x^n}{n} = x + \frac{x^2}{2} + \frac{x^3}{3} + ...\]
  converges for every $x$ in the half-open interval $[-1, 1)$ but does not converge when $x=1$. For a fixed $x_{0} \in (-1, 1)$, explain how we can
  still use the Weierstrass M-Test to prove that $x$ is continuous at $x_{0}$.
\end{enumerate}

\begin{solution}
  \begin{enumerate}[label=(\alph*)]
    \item We have for $x \in [-1, 1]$
    \[\left|\frac{x^n}{n^2}\right| \leq \frac{1}{n^2}\]
    and $\sum_{n=1}^{\infty}\frac{1}{n^2}$ converges. This satisfies the conditions for the Weierstrass M-Test, hence the series converges uniformly.

    We have $\frac{x^n}{n^2}$ continuous for every $n$. Combined with the uniform convergence above, we may apply
    the Term-by-term Continuity Theorem to conclude that $h$ is continuous.

    \item For any $x_{0} \in (-1, 1)$, we can find an $r$ such that $|x_{0}| < r < 1$. Now we can bound the function with $r$. For every $x \in [-r, r]$:
    \[\frac{x^{n}}{n} \leq \frac{|x^n|}{n} \leq r^{n}\]
    Since $|r| < 1$, the geometric series $\sum_{n=1}^{\infty}r^{n}$ converges and we can apply the Weierstrass M-Test to conclude that
    $\sum_{n=1}^{\infty}\frac{x^n}{n}$ converges uniformly on $[-r, r]$.

    Since each term of the series is continuous, we may apply the Term-by-term Continuity Theorem to conclude that
    $f(x)$ is continuous on $[-r, r]$. Since $x_{0} \in [-r, r]$, $f$ is continuous on $x_{0}$ as well.
  \end{enumerate}
\end{solution}

%------------------------- Problem 6 -------------------------
\miquestion Let
\[f(x) = \frac{1}{x} - \frac{1}{x+1} + \frac{1}{x+2} - \frac{1}{x+3} + \frac{1}{x+4} - ...\]
Show $f$ is defined for all $x > 0$. Is $f$ continuous on $(0, \infty)$? How about differentiable?

\begin{solution}
  By the Alternating Series Test, the series converges for all $x > 0$.
  
  It's continuous on $(0, \infty)$. We use the same trick as problem 4 and 5. For any $x \in (0, \infty)$, there exists $a, b$ such that $0 < a \leq x \leq b$.
  We'll prove that $f$ converges uniformly on $[a, b]$. Then by the Term-by-term Continuity Theorem, $f$ is continuous on $[a, b]$.

  We'll prove uniform convergence by using the Cauchy Criterion (this is very similar to Exercise 2.7.1).
  First, recall the following fact about partial sums of an Alternating Series
  \[|s_j - s_{i}| = |\sum_{k=i+1}^{j}(-1)^{k+1}a_{k}| \leq a_{i+1}\]

  Pick $N$ large enough such that $\frac{1}{a+N} < \epsilon$. Note that because the domain is $[a, b]$, then for any $x \in [a, b]$
  \[\frac{1}{x+N} \leq \frac{1}{a+N} < \epsilon\] 
  Now for any $n > m > N$ and any $x \in [a, b]$, we have
  \[|s_{n}(x)-s_{m}(x)| \leq \frac{1}{x+N} \leq \frac{1}{a+N} < \epsilon\]
  satisfying the Cauchy Criterion.

  It's differentiable. We'll show that it satisfies the conditions of the Term-by-term Differentiability Theorem.
  
  Condition \#1: We have $f_{n}(x)=\frac{(-1)^n}{x+n}$ is differentiable with $f_{n}'(x)=-\frac{(-1)^{n+1}}{(x+n)^2}$.

  Condition \#2: $\sum_{n=1}^{\infty}f_{n}'(x)$ converges uniformly by the Weierstrass M-test. The bound is
  \[\left|\frac{(-1)^{n+1}}{(x+n)^2}\right| \leq \frac{1}{n^2}\]
  and $\sum_{n=1}^{\infty}\frac{1}{n^2}$ converges.

  Condition \#3: We've already established that $f$ is defined for all $x > 0$ (the Theorem only requires a single point $x_{0} \in [a, b]$ where $f$ is defined).
  Now all conditions are satisfied.
\end{solution}

\end{questions}
\end{document}
